\subsection*{先给事件命名，再讨论先后}

Alice 的客户端、购买服务和数据库各自执行自己的步骤，也各自取得消息。先给几个关键事件命名：

\begin{center}
\begin{minipage}{\linewidth}
\small
\begin{tabularx}{\linewidth}{p{.12\linewidth}X}
\toprule
\textbf{事件} & \textbf{含义} \\
\midrule
C1 & 客户端发出请求 \\
S1 & 服务受理请求 \\
D1c & 数据库完成本次提交 \\
S2 & 服务收到提交成功确认 \\
S3 & 服务发出 HTTP 成功响应 \\
C2 & 客户端收到并接受完整业务结果 \\
CT & 客户端等候结果的截止时间到达 \\
\bottomrule
\end{tabularx}
\end{minipage}
\end{center}

这里用 D1c 表示提交事件，扣款流水仍叫 D1。一条执行历史，是把某次运行中发生的发送、处理、提交等事件及其先后关系整理在一起。下面的箭头表示因果先后：右边的事件沿这条处理路径发生在左边之后。在一条正常成功历史中，有如下因果链：
\[
  \mathrm{C1}\longrightarrow\mathrm{S1}\longrightarrow\mathrm{D1c}
  \longrightarrow\mathrm{S2}\longrightarrow\mathrm{S3}\longrightarrow\mathrm{C2}.
\]

服务先受理请求，随后才开始业务处理；按当前规则，服务收到提交成功确认后才发送成功响应。每个进程内部的程序顺序、一次消息的发送与接收，以及这些关系的连续传递，给出了这里的事件顺序。

CT 是另一个本地事件。客户端只需依据自己的等待状态和时钟，就能结束等待。数据库可能还没接到事务，也可能已经提交、只是响应尚未到达。这两种情况都允许 CT 发生。

\begin{figure}[htbp]
\centering
\begin{tikzpicture}[x=1cm,y=1cm]
\foreach \x in {0,4.1,8.4,12.5} {\draw[Ink!35,dashed] (\x,-.6)--(\x,-4.7);}
\node[font=\small\sffamily\bfseries] at (2.05,.7) {提交先完成，响应迟到};
\node[font=\small\sffamily\bfseries] at (10.45,.7) {客户端先超时，提交稍后完成};
\node[booknode,text width=1.5cm] at (0,0) {客户端};
\node[booknode,text width=1.5cm] at (4.1,0) {数据库};
\node[booknode,text width=1.5cm] at (8.4,0) {客户端};
\node[booknode,text width=1.5cm] at (12.5,0) {数据库};
\draw[bookarrow] (0,-1.1)--node[above,font=\scriptsize]{请求与执行路径}(4.1,-1.7);
\node[booknode,text width=1.7cm] at (4.1,-2.4) {D1c：提交};
\node[bookresult,text width=1.7cm] at (0,-3.6) {CT：超时};
\draw[bookarrow] (8.4,-1.1)--node[above,font=\scriptsize]{请求与执行路径}(12.5,-1.7);
\node[bookresult,text width=1.7cm] at (8.4,-2.4) {CT：超时};
\node[booknode,text width=1.7cm] at (12.5,-3.6) {D1c：提交};
\node[booknote,text width=11cm] at (6.25,-5.2) {两条历史中，客户端都在截止时间前没有取得完整业务响应};
\end{tikzpicture}
\caption{本地超时可以出现在不同远端进度下。请求和消息建立因果关系，截止时间由客户端独立触发。}
\label{fig:ch02-causality}
\end{figure}


服务收到本次 COMMIT 成功确认，然后才暂停响应路径，这条程序路径就把提交确认排在了故障之前。这里设置的暂停点也称故障屏障：它让执行停在指定位置，便于观察随后发生的超时。恢复后再读到 P1 和关联记录，可以核对最终保留的业务事实。一条准备提交的日志则只记录进程走到准备提交的位置，还需要提交确认或已提交记录说明后来的结果。

\subsection*{三个参与者分别知道什么}

客户端在发出 R1 后，知道自己试图购买 \code{paper-001}，也知道请求使用 Alice 的凭据。凭据是服务验证调用者身份的材料，本例请求头中的访问令牌用于让服务认出 Alice。若客户端只收到 TCP 层的确认，尚未收到业务响应，P1 对它仍然未知。TCP 在连接两端传送有序字节，购买请求和响应由这些字节承载。CT 到达以后，客户端新取得的信息是本地等待期限已经耗尽。

服务知道自己受理了 R1，并生成了 P1。条件 \code{UPDATE} 返回一行时，扣款步骤已经在当前事务内完成；后续错误仍可能让整笔事务回滚。服务收到 COMMIT 的成功确认后，才按当前数据库配置和故障范围，把 P1 作为已提交购买返回。

数据库处理事务和记录，通常不知道 Alice 是否看到了页面。R2 到来时，它也缺少信息判断这代表第二次独立购买，还是继续尝试第一次购买。应用提交了哪些字段和关系，数据库约束就据此检查哪些关系。

\begin{center}
\begin{minipage}{\linewidth}
\centering\small\captionsetup{hypcap=false}
\captionof{table}{每项观察提供的证据}
\begin{tabularx}{\linewidth}{>{\raggedright\arraybackslash}p{.24\linewidth}XX}
\toprule
\textbf{位置与事件} & \textbf{已经取得的证据} & \textbf{仍需另行取得的信息} \\
\midrule
客户端发送完成 & 请求已交给本地发送路径 & 服务执行与事务提交结果 \\
服务收到条件更新结果 & 当前事务内的余额更新结果 & 整笔事务的最终提交结果 \\
服务收到提交成功确认 & 本事务已按数据库保证提交 & 客户端取得业务结果的情况 \\
客户端接受完整成功响应 & 依照接口规则确认购买成功 & 以后查询所需历史的保留情况 \\
恢复者读到购买终态和关联记录 & 数据库保留了这笔购买的结果 & 原客户端显示过什么 \\
\bottomrule
\end{tabularx}
\end{minipage}
\end{center}

客户端通过成功响应判断购买完成，依赖几个具体条件：响应来自可信服务，属于正在处理的调用，内容完整，服务执行先提交再报告成功的规则。在这些条件下，响应把已提交购买的信息传到了客户端。

数据库的持久性也有明确前提。持久性指事务提交以后，记录在约定的故障和恢复范围内仍然保存。本章采用保留已确认提交的数据库配置，服务重启后仍连接这份数据库。存储设备失效或恢复旧备份可能让已确认记录丢失，需要另行讨论数据恢复和历史回退后的处理规则。

\subsection*{TCP ACK、HTTP 响应和 COMMIT 确认}

同一笔购买会出现几种确认。TCP ACK 是传输层对已接收字节的确认。TCP 给连接中的字节编号，确认号表示对端 TCP 已经收到该序号以前的连续字节。它用于处理字节传输、重传和顺序，对端应用这时可能尚未读取或校验请求，也可能尚未开始事务。

服务返回 HTTP 响应，说明接口层的处理结果。服务先完成提交，再发出成功；客户端完整收到并识别这份结果，才获知购买已经完成。连接中途关闭、响应只到了一部分，或者响应和待处理调用的关联丢失，都会使这个信息传递过程停在中间。响应中的三位状态码概括本次接口处理情况。本例用 \code{503} 表示服务暂时无法处理，用 \code{409} 配合业务字段报告余额不足\footnote{503 的标准含义为 Service Unavailable，即服务不可用；409 为 Conflict，即请求与资源当前状态冲突。将余额不足映射为 409 是本例的接口选择，具体原因仍由业务字段说明。本书中的 201 表示创建成功，200 表示请求成功，404 表示未找到；各状态码都须结合相应入口的规则理解。}。

COMMIT 的确认发生在服务与数据库之间。服务发出 COMMIT 后，需要检查数据库协议中报告命令完成的标签、错误信息与事务状态。完成标签说明刚才哪条命令完成；提交成功要有与 COMMIT 对应的成功信息。PostgreSQL 的 \code{ReadyForQuery} 是告知客户端本轮查询字符串已处理完、可以接收下一轮查询的协议消息，出错以后也会出现。连接恢复可用与本次提交成功，是两项不同的信息。

数据库可能已提交，提交确认却在返回途中丢失。服务此时只看到数据库连接异常，面对的也就成了一项尚未获知结果的操作。

\begin{figure}[htbp]
\centering
\begin{tikzpicture}[x=1cm,y=1cm]
\node[booknode,text width=2.3cm] (a) at (0,0) {发送端 TCP};
\node[booknode,text width=2.3cm] (b) at (9.5,0) {接收端 TCP};
\draw[bookarrow] (b)--node[above,font=\small]{ACK：对应序号之前的字节已接收}(a);
\node[booknote] at (4.75,-.6) {应用读取、业务校验和数据库处理仍在后续路径中};
\node[booknode,text width=2.3cm] (c) at (0,-1.8) {购买服务};
\node[booknode,text width=2.3cm] (d) at (9.5,-1.8) {数据库};
\draw[bookarrow] (d)--node[above,font=\small]{COMMIT 成功确认：本事务已提交}(c);
\node[booknote] at (4.75,-2.4) {服务取得提交结果，客户端还在等待};
\node[booknode,text width=2.3cm] (e) at (0,-3.6) {客户端};
\node[booknode,text width=2.3cm] (f) at (9.5,-3.6) {购买服务};
\draw[bookarrow] (f)--node[above,font=\small]{完整业务响应：P1 成功，O1、D1}(e);
\node[booknote] at (4.75,-4.2) {客户端依据服务先提交再返回的规则，获知购买结果};
\end{tikzpicture}
\caption{三种确认分别说明字节接收、事务提交和业务结果。每一层确认都有自己的对象。}
\label{fig:ch02-ack-layers}
\end{figure}


TCP 处理传输中的重复数据包；应用重建连接，再发送整个 POST，新字节流可以承载一笔新的业务请求。描述 TCP 的标准文档 RFC 9293 规定了可靠有序字节流和累计确认：已经确认到某个序号，就同时确认它以前的连续字节\cite{tcp}。RFC 是 Request for Comments 的缩写，这里按编号指向具体的互联网协议规范。TCP 的保证范围止于对应的传输行为；R1 与 R2 的购买关系仍由应用表达。

描述 HTTP 语义的 RFC 9110 第 9.2.2 节按用户请求的预期效果定义幂等方法：重复发送相同请求时，对服务端产生的预期效果与发送一次相同。该节还讨论连接失败以后的重试。对于非幂等方法，规范使用 SHOULD NOT 限制自动重试。SHOULD NOT 是规范中的要求用语，含义是通常不应这样做，只有理解相关后果并有充分理由时才考虑例外。这里的例外包括客户端已知请求的实际语义可安全重复，或者能够发现原请求从未应用\cite{http}。当前入口每次 POST 都创建新购买，重发会产生新效果。

增加 TCP 重传次数、延长 HTTP 等待时间或调整数据库连接超时，会改变客户端多久遇到未知结果。它们没有在购买数据中增加 R1 与 R2 的关系。

\subsection*{成功提交与获知提交是两个事件}

数据库提交 P1 以后，满足读取条件的后续事务可以看见它。服务需要消息或查询获知这个事实，客户端还需要更远的一次消息传递。

在提交后响应丢失的历史里，数据库保存了 P1，服务曾收到提交确认，客户端尚未获知结果。若故障发生在数据库确认返回以前，数据库可能保存了 P1，服务和客户端都还未知。若客户端收到成功后、保存 P1 前退出，重启后的客户端又可能失去查询依据。

\begin{figure}[htbp]
\centering
\begin{tikzpicture}[x=1cm,y=1cm]
\node[booknode,text width=2.7cm] (db) at (0,0) {数据库保存\\P1、O1、D1};
\node[booknode,text width=2.7cm] (svc) at (4.8,0) {服务收到\\提交成功确认};
\node[bookresult,text width=2.7cm] (cli) at (9.6,0) {客户端等待\\完整业务结果};
\draw[bookarrow] (db)--node[above,font=\scriptsize]{确认到达}(svc);
\draw[bookalert,dashed] (svc)--node[above,font=\scriptsize]{响应丢失}(cli);
\node[booknode,text width=2.7cm] (disk) at (0,-2) {重启后\\仍可读取记录};
\node[booknode,text width=2.7cm] (recovery) at (4.8,-2) {恢复进程\\重新查询 P1};
\node[bookresult,text width=2.7cm] (lost) at (9.6,-2) {客户端退出\\内存中的 P1\\可能丢失};
\draw[bookarrow] (db)--(disk);
\draw[bookarrow] (disk)--node[above,font=\scriptsize]{持久事实}(recovery);
\draw[bookalert] (cli)--(lost);
\node[booknote,text width=11.8cm] at (4.8,-3.3) {恢复要接通两端：数据库保留结果，调用方保留能定位原操作的信息};
\end{tikzpicture}
\caption{事实保存、服务获知、客户端获知分别发生。持久记录与客户端的恢复依据需要接续起来。}
\label{fig:ch02-fact-propagation}
\end{figure}


Birrell 与 Nelson 在 1984 年的《Implementing Remote Procedure Calls》中，设计了把远程通信组织为过程调用的 RPC 系统。RPC 即远程过程调用：调用方给出过程名和参数，本地通信代码把它们送到另一台机器执行，再把结果带回。论文中的 RPC 运行时负责重传和识别重复调用，并区分正常返回与通信异常的含义：调用正常返回时，服务器过程被调用恰好一次；发生所述异常时，过程可能调用过一次，也可能完全没有调用，调用者没有获知是哪一种\cite{rpc}\footnote{原文对异常一侧的说明为：“the procedure will have been invoked either once or not at all---the user is not told which.”}。该论文讨论的是它自己的 RPC 运行时。过程被调用一次，与本书的数据库事务提交一次，仍各有对应的处理步骤。

购买客户端收到超时，也处在这样的信息缺口里。它知道这次通信没有给出完整结果；远端的业务结果则需要由响应、查询或恢复过程继续查明。页面可以显示待确认，并保留继续处理所需的信息。

\subsection*{日志怎样帮助复原历史}

排查这笔购买时，需要把客户端的观察、服务的执行位置和数据库留下的记录联系起来。R1、P1、O1、D1 的关联日志帮助找到相应记录，数据库关系则用于核对业务效果。

日志停在发出 COMMIT，说明服务走到了发送提交命令的位置；后续结果仍需从提交确认或数据库中查明。日志停在准备响应，说明服务开始组织返回内容；客户端实际收到了什么，还需要发送路径和客户端记录。

若服务记录了 \code{request\_id=R1}、\code{purchase\_id=P1} 和 \code{commit\_ack\_received}，数据库中也有 \code{succeeded} 的 P1 及完整 O1、D1，就能确认本次购买已经提交。响应具体在哪里中断，则是接下来要核对的链路。

成功 HTTP 状态码的数量与扣款数量可能不同。R1 没有一条完成的访问日志，数据库却已有一次扣款；R2 返回了一次 201，账本中实际有两条流水。审计需要逐笔核对流水，再通过编号追溯相应尝试。
\FloatBarrier
